\chapter*{Appendix B: User Manual}
\addcontentsline{toc}{chapter}{Appendix B}
\markboth{Appendix B}{Appendix B}

\section*{Prerequisites}

The fuzzer requires Python 3.10+ and the \texttt{uv} package manager. For the Rust echo server, a Rust toolchain is needed. TLS certificates (self-signed) are used for local testing.

\section*{Basic Usage}

\begin{lstlisting}[language=bash,frame=single]
# Validate connection
uv run python main.py --no-fuzz \
    --url https://127.0.0.1:6161/echo

# One-shot capsule fuzzer
uv run python main.py --mode oneshot \
    --url https://127.0.0.1:6161/echo

# Multistep scenario fuzzer
uv run python main.py --mode multistep \
    --url https://127.0.0.1:6161/echo

# With server subprocess management
uv run python main.py --mode multistep \
    --url https://127.0.0.1:6161/echo \
    --server-cmd "uv run server/aioquic/server.py \
        cert.pem key.pem 2>&1 | tee server.log"

# Resume from a specific test case index
uv run python main.py --mode multistep \
    --start-index 42 --end-index 100
\end{lstlisting}

\section*{Server Fingerprinting}

Before fuzzing, identify which draft generation the target implements:

\begin{lstlisting}[language=bash,frame=single]
uv run get-server-version.py --no-verify \
    127.0.0.1 6161
\end{lstlisting}

\section*{Analyzing Results}

After a fuzzing run, correlate test cases with server output:

\begin{lstlisting}[language=bash,frame=single]
uv run analyze_logs.py --log server.log \
    --db boofuzz-results/run_<timestamp>.db
\end{lstlisting}

The BooFuzz web UI is available at \texttt{http://localhost:26000} during a run.
